% TOP_PROBLEM: {"id": 5200013, "title": "Open Problems on Billiards and Geometric Optics", "category_id": 11, "category": {"id": 11, "name": "dynamical_systems", "display_name": "Dynamical Systems", "description": "Problems about long-term behavior of deterministic systems, Hamiltonian dynamics, and periodic orbits.", "slug": "dynamical-systems", "order_index": 11, "created_at": "2026-07-31T15:26:25.671Z"}, "status": "partially_solved"}
% FIRST_POSTED: null
% ENTRY_KIND: statement_only
% Imported from: batch03.tex; UnsolvedMath ID 5200013
\documentclass[11pt]{article}
\usepackage[margin=1in]{geometry}
\usepackage{amsmath,amssymb,mathtools}
\usepackage{fontspec,unicode-math}
\setmainfont{Latin Modern Roman}
\setsansfont{Latin Modern Sans}
\setmathfont{Latin Modern Math}
\usepackage{xurl,hyperref}
\hypersetup{colorlinks=true,urlcolor=blue,linkcolor=blue}
\newcommand{\sref}[2]{\hyperlink{source:#1:#2}{[S#2]}}
\newcommand{\cataloguescope}[1]{\paragraph{Recorded mathematical question.}#1}
\begin{document}
% UnsolvedMath ID 5200013; source catalogue code AMR-051-0013
\setcounter{section}{5200012}
\section{Open Problems on Billiards and Geometric Optics}\label{5200013}
{\small\sffamily UnsolvedMath ID 5200013 · Dynamical Systems}
\subsection{Definitions and mathematical statement}

\paragraph{Shared notation.}$\mathbb N=\{1,2,\ldots\}$, and $\mathbb Z,\mathbb Q,\mathbb R,\mathbb C$ denote the integers, rationals, reals and complex numbers. Any other conventions are specified locally.


A ray is an oriented Euclidean line before entering a mirror system. Trapping means that after entry its successive optical reflections keep it in the trapping domain and prevent escape. The line space $\mathcal L$ has dimension $2n-2$ and its standard symplectic form. A local normal family consists of rays perpendicular to a hypersurface; it is a Lagrangian submanifold of line space of dimension $n-1$. A non-normal family lacks this property.
\paragraph{Statement recorded in the supplied source.}
In $\mathbb{R}^n$, the space $\mathcal{L}$ of oriented lines has dimension $2n-2$ and a natural symplectic structure. Normal families of rays form Lagrangian submanifolds and can be trapped by mirrors. What is the greatest dimension of a family of rays that can be trapped? In $\mathbb{R}^3$, can a non-normal two-parameter family of rays be trapped? \sref{5200013}{2}
\subsection{Short English statement}
Find the maximum dimension of a family of rays that mirrors can trap in Euclidean space; in three dimensions, decide whether a non-normal two-parameter family can be trapped.
\subsection{Sources}
\begin{description}
\item[\hypertarget{source:5200013:1}{S1}] ULAM.AI and the UnsolvedMath Contributors, \emph{UnsolvedMath: A Curated Collection of Open Mathematics Problems}, record 5200013 (AMR-051-0013). CC BY 4.0. \url{https://huggingface.co/datasets/ulamai/UnsolvedMath}.
\item[\hypertarget{source:5200013:2}{S2}] M. Bialy, C. Fierobe, A. Glutsyuk, M. Levi, A. Plakhov and S. Tabachnikov, \emph{Open problems on billiards and geometric optics} (2021), arXiv:2110.10750v2, Serge Tabachnikov, Problem 1, pp. 9--10. \url{https://arxiv.org/pdf/2110.10750v2}.
\end{description}
\label{end:5200013}
\end{document}
